% Part III introduction (no \chapter heading; flows after the \part page).

\begin{draftblock}
\noindent The mechanisms of \cref{part:oversight} presuppose that we can elicit honest beliefs from agents. Economics has a canonical answer for how to do this: \emph{proper scoring rules} \citep{goodRationalDecisions1952, savageElicitationPersonalProbabilities1971, gneitingStrictlyProperScoring2007} reward a forecast against the realized outcome so that truthful reporting maximizes expected reward, and \emph{prediction markets} \citep{hansonLogarithmicMarketScoring2002, hansonCombinatorialInformationMarket2003} decentralize this: an asset contingent on the event is traded, and an agent's ``fair price'' is its subjective probability. Indeed the two are equivalent in a precise sense: a logarithmic market scoring rule is exactly a sequentially-shared proper scoring rule, and the market subsidy is the price paid for information.

Two gaps separate this classical toolkit from what alignment needs.

\textbf{First, the questions we care about may never resolve.} Scoring rules and prediction markets presuppose a \emph{verifiable or falsifiable} (VF) event: ground truth eventually arrives and bets are settled. But human values, latent variables, counterfactuals and ``subjective'' questions are not VF --- there is nothing to settle the bet against. A classical observation in this direction is due to \citet{conitzerPredictionMarketsMechanism2012}, who studied prediction markets through the lens of mechanism design and cooperative game theory, and proposed a \emph{two-layer} structure: an object-level market that prices claims against ground truth, and a second layer that distributes credit among informants according to their marginal contribution to the market's prediction --- so that agents can be paid for \emph{information} even when their individual contribution is never directly verified, so long as the \emph{bottom} layer is eventually anchored to ground truth. The question this part pursues is what happens when even that anchor is removed.
%% TODO(part3-intro): elaborate exactly how the subjective-scoring level connects to
%% economics (per thesis plan), and verify the characterization of Conitzer (2012)'s
%% two-layer construction against the paper before submission.

\textbf{Second, even where ground truth exists, it may be too slow.} Evaluating (and training) AI forecasters against resolution alone means waiting months or years for feedback. But rationality constraints --- coherence of probabilities across logically related questions --- can be checked \emph{immediately}, and violations of them are exactly \emph{arbitrage opportunities} in the market reading of beliefs.
\end{draftblock}

% Adapted from the year-2 report (P4, P5):
For the special case where ``non-VF'' sentences may be represented in a first-order logic, \cref{ch:nonvf} develops a mechanism for betting on them called the \emph{verification-falsification game}: the key idea being that although such sentences never resolve, agents can be made to bet on them through a game whose \emph{subgames} concern VF instances quantified over by the original sentence. One way to evaluate forecasts on sentences that are non-VF -- or more generally, might take a lot of time or effort to get ground truth on -- is to use \emph{consistency checks}: \cref{ch:consistency} develops a consistency benchmark for evaluating AI forecasters, in particular a metric based on market arbitrage to measure inconsistency, and a market-based method to correct inconsistency analogous to what Platt scaling is for miscalibration.
